% Self-contained comparison section.  Requires amsmath, booktabs, longtable,
% xurl (or url), and hyperref.  No project-specific macro is used.
\section{Comparison of the supplied constructions}
\label{sec:comparison-audit}

\subsection{Scope, model, and exact source identities}
All state counts in this comparison refer to working states of a deterministic
binary Turing machine, with an additional halting label, a tape infinite in
both directions, and an initially blank tape.  A missing working-state number
is not a state: the census counts distinct source labels in the transition
table, not the largest label plus one.  In particular the historical
PA209 table has 209 working labels and halt label 215, whereas PA196 uses the
dense working labels $0,\ldots,195$ and halt label 196.

The comparison fixes the following repository revisions.  The links name
immutable revisions; references to a repository's moving \texttt{main} branch
are not used to identify a table.
\begin{description}
\item[ProveIt incoming and ConciseIPRA.]
\href{https://github.com/VladimirReshetnikov/ProveIt/tree/6ce5acc401f0c84be90dc5c38dc1487e391122c6}{ProveIt revision}
\nolinkurl{6ce5acc401f0c84be90dc5c38dc1487e391122c6}.
\item[Published PA196.]
\href{https://github.com/jamesyc/BusyBeaver/tree/407cdc356f6b586635e71ac92be394255f4bfea4}{BusyBeaver revision}
\nolinkurl{407cdc356f6b586635e71ac92be394255f4bfea4}.
\item[Historical PA209.]
\href{https://github.com/jamesyc/BusyBeaver/tree/69d98241a721ead20b4f0b17df9062950b40fbfe}{BusyBeaver revision}
\nolinkurl{69d98241a721ead20b4f0b17df9062950b40fbfe}.
\end{description}

The complete literal-file SHA-256 digests are:
\begin{description}
\item[\texttt{data/pa196.tm}]
\nolinkurl{e8bbccd4cfa325251ca95066db116b5aa766e404b0971c33c02772be1c3617a8}.
\item[\texttt{data/pa209.tm}]
\nolinkurl{0610d978675320f0b4b1ef34619dcd3daa5912632d7037e200dfcc70954fea51}.
\end{description}
These are hashes of the supplied files, including their literal row order and
formatting.  They must be distinguished from the hashes of canonical
transition graphs used below.  The audit manifest records both kinds of
identifier and the complete provenance of every selected table.

\subsection{The incoming archive census: twenty graph classes}
All eleven ZIP archives in the pinned \texttt{docs/incoming} directory were
inventoried.  Their names do not denote eleven independent reductions.  Some
contain alternate presentations of the same table, some contain further
208-state tables, and others record negative searches or uncompiled logical
proposals.  Table~\ref{tab:incoming-census} separates these roles.  Except for
\texttt{pa\_unary\_power\_cover}, the archive names have the suffix
\texttt{\_2026-10-29.zip}; the shorter stems are printed for readability.

\begin{longtable}{@{}p{.39\linewidth}p{.57\linewidth}@{}}
\caption{The eleven supplied incoming archives.}\label{tab:incoming-census}\\
\toprule
Archive stem & Audited role\\
\midrule
\endfirsthead
\toprule Archive stem & Audited role\\\midrule
\endhead
\path{U208_ConPA} & Primary 208-state presentation in original graph class A.\\
\path{U208_ConPA_Conditional_Context} & Primary A and alternative B; the retained member of the folded pair differs from the first presentation.\\
\path{U208_ConPA_Entry_Context} & Primary A and alternative B; its primary and alternative literal tables duplicate the corresponding Conditional Context files.\\
\path{U208_ConPA_History} & Primary A and alternative B, with a history-based explanation of the same local reduction.\\
\path{u208_research} & Primary B and a Plain alternative in A.\\
\path{pa208_research} & Eighteen additional 208-state layouts: nine modifications of each original graph class.\\
\path{Below_U208_Research} & Boundary projections and guarded, offset, and source searches; no accepted smaller table.\\
\path{ConPA_below208_research} & Comparison, two-row searches, source permutations, and a finite-model obstruction to a logical simplification.\\
\path{pa_consistency_research} & Source specialization from 3,909 to 3,904 primitives, but an increase from 197 to 200 syntactic subtrees; no smaller compiled table.\\
\path{u208_comparison_research} & Comparison and local/source searches; the expressly rejected 207-state candidate is not an accepted consistency machine.\\
\path{pa_unary_power_cover} & A mathematical unary-cover proposal and a separate proof stream, without a compiled literal Turing machine.\\
\bottomrule
\end{longtable}

The first five rows supply nine designated original 208-state presentations.
A fresh, independent finite verifier classifies these into exactly two rooted,
action-preserving transition-graph classes, denoted A and B.  The additional
eighteen designated tables in \texttt{pa208\_research} belong to eighteen more
classes.  Thus the selected accepted presentations give
\[
  9+18=27\text{ presentations},\qquad 2+18=20\text{ graph classes}.
\]
Here ``accepted'' means the archive's designated final 208-state
construction, together with the independent finite local checks described
below.  It does not reclassify failed exploratory tables as valid machines,
or assert a new proof-assistant certification of the original packages.

The classification does more than compare file hashes.  Starting at state 0,
the verifier numbers working states by breadth-first discovery, following
read-zero before read-one.  It retains the read symbol, written symbol,
head direction, and target, with a separate distinguished halt symbol.
Every working state is graph reachable.  Equal serializations therefore give
an exact rooted action-preserving isomorphism, with the bijection explicitly
available from the traversal.  Different serializations prove that no such
isomorphism exists.  The two original canonical SHA-256 identifiers, in this
audit's specified serialization, are
\begin{description}
\item[A]
\nolinkurl{43d44b9b61daf753b2b833273c6b970c1787b15b7e4e177a30ca4cd886486702}.
\item[B]
\nolinkurl{7bd9c53560bd49e008cdf4653800e0218c5ec94330e24d2a714f35c26971d63b}.
\end{description}
Canonical hashes produced by a different serialization need not agree.
The complete file identities and the eighteen other canonical identifiers
are retained in \texttt{inputs/manifest.json} and
\texttt{comparison\_results.json}.

\subsection{The original one-state U208 reduction}
The five original 208-state packages all exploit the same pair of historical
native states, 172 and 211.  Their savings cannot be added.  In the
Joint/History naming, retain 172, delete 211, redirect both outgoing rows of
163 that previously entered 211 to 172, and replace the read-zero action of
172 by
\[
               172,0\longmapsto(0,L,217).
\]
An equivalent naming keeps 211 and deletes 172, redirecting the appropriate
entry from 116.  The differing retained label is not a second deletion.

The justification depends on how the state is entered.  Complete predecessor
histories establish that an old entry to 172 has the two preceding tape bits
$10$, whereas an old entry to 211 that scans zero has preceding bit $1$.
The retained state can move left and use that history bit to select the old
behavior.  The exceptional transitions expand to exact finite blocks of
lengths $7,1,3,1$.  At each block endpoint the entire tape and head position
agree with the old endpoint, with only the state identification applied.
All other transitions commute with that identification.

For an aligned finite old execution, let $N(q,b)$ count old steps in state
$q$ scanning $b$.  The construction gives
\[
 T_{208}=T_{209}+6N(172,0)+2N(211,0),\qquad
 T_{209}\leq T_{208}\leq7T_{209}.
\]
These are positive blocks, so the correspondence preserves both halting and
nonhalting.  The induction starts in the specified initial control state and
applies to every common initial binary tape; it is not based on a long
blank-tape sample run.

The two original U208 source layouts have the same ordered 3,909 primitive
opcodes and operands and the same 5,083 ordinal control-flow edges.  Their
binary associations and physical dispatchers differ.  The original packages
record an exact startup of 48,140 steps and 5,084 uniform source-site
contracts.  This comparison inspected those proof artifacts but did not rerun
their complete older symbolic-runtime derivations.  The graph census and
the extra local helper verifications are independently reproduced here.

\subsection{The eighteen additional 208-state layouts}
The additional tables reuse a related idea without reducing the final state
count.  They delete one state $q$, retain a state $r$, and add a new helper,
named 500 in the native tables.  On the exceptional scanned bit $b$, the
retained state moves left to read the tag left by its predecessor.  The
helper dispatches to existing continuations.  Table~\ref{tab:helper-families}
lists all choices; each is applied once to each original graph class.

\begin{table}[htbp]
\centering\small
\begin{tabular}{@{}rrcllc@{}}
\toprule
Deleted $q$&Retained $r$&$b$&Helper on 0&Helper on 1&Choices\\
\midrule
127&124&0&$0R8$&$1R59$, $1R109$, or $1R125$&3\\
190&160&1&$0R86$&$1R193$&1\\
178&168&0&$0R154$ or $0R159$&$1R176$ or $1R225$&4\\
301&88&0&$0R122$&$1R500$&1\\
\bottomrule
\end{tabular}
\caption{Nine neutral helper layouts over each of A and B.  All labels are
historical native labels, not PA196 dense labels.}
\label{tab:helper-families}
\end{table}

Each layout has a net change of zero working states: one deletion and one
new helper.  The independent verifier checks the exact literal patch, the
complete incoming-edge census, and the fixed opposite predecessor bits.
It then checks the four relevant two-cell executions while treating the
exterior tape as opaque.  Every macro has positive length at most five,
preserves the entire exterior, and ends at the required old endpoint.
Ordinary rows satisfy the quotient correspondence.  These finite checks
establish the local simulation for all allowed surrounding tapes, rather
than just sampled register values.  They also explain why merely noticing
that two rows have similar continuations is insufficient: the predecessor
tag and all uses of the chosen helper must be compatible.

\subsection{How published PA209 and PA196 differ from the U208 family}
The U209 parents used inside the ProveIt U208 archives are not literally the
published PA209 machine at the historical BusyBeaver revision.  The shared
engine and compiler lineage explain many common native labels and rows,
but the proof-command bodies, source layouts, and startups differ.  For
example, History U208 class A and published PA209 have 195 common native
working labels and 304 identical transition rows.  This is a useful
comparison in the inherited naming, not a conjugacy claim.

The principal quantitative comparison is
\begin{center}
\begin{tabular}{@{}lrrr@{}}
\toprule
Construction&Working states&Physical source sites&Blank startup steps\\
\midrule
Original U208 classes A/B&208&3,909&48,140\\
Published PA209&209&3,948&58,210\\
Published PA196&196&3,812&58,154\\
\bottomrule
\end{tabular}
\end{center}
The PA209 source-site count is stated by the \texttt{sites\_count}
theorem in its pinned \texttt{M6/Program.lean}; the PA196 source-site count is
stated in \texttt{PA196/Program.lean}.  This audit did not rerun those upstream
Lean builds.  The published PA196 paper describes its physical
state partition as 36 engine states, 13 leaf routes, and 147 states reading
159 distinct source-tree nodes.

PA196 obtains its thirteen-state improvement by changing the logical command
construction and its compiled layout together.  Its universal-quantifier
constructor transposes the Cantor-pairing arguments:
\[
                    U(a,B)=4p(B,a)+2.
\]
It deletes the old primitive equality/transitivity head $\rho_7$, adds guarded
equality vacuity to the gate, and derives the required reverse membership
rewrites and neighborhood equality using separation.  It also mirrors the
equality witness, exchanges selected antecedents and axiom arguments,
reorders the command items, and stores the resulting 201-item AX body as
68 internal nodes represented by 62 physical states through six
context-dependent pairs.  Two obsolete outer controls are removed.
These are described in the pinned
\href{https://github.com/jamesyc/BusyBeaver/blob/407cdc356f6b586635e71ac92be394255f4bfea4/paper/pa196.tex}{PA196 paper}
and its associated command and source modules.

An exact native-label table comparison removes 75 old working labels,
introduces 62 new labels in the 1300 series, and retains 134 working labels.
Of the retained rows, 265 are identical and only the following three change:
\[
\begin{array}{cc|l|l}
q&b&\text{old action}&\text{new action}\\\hline
101&1&1R115&1R1367\\
164&0&0L311&0R77\\
190&1&1L308&1R99.
\end{array}
\]
The net census is $209-75+62=196$.  In particular both native 172 and native
211 are already absent from PA196.  The central U208 deletion therefore
cannot be applied again to remove one of them.

Two neutral U208 helper families, 127/124 and 301/88, survive far enough
structurally to motivate new searches.  However the old helpers cannot simply
be overlaid on PA196.  The independently attempted existing-helper
substitutions for 127/124 fail uniform source-site contracts, and the
obvious self-loop donors for 301/88 belong to the protected register engine.
More extensive finite searches are recorded separately with their precise
scope and counterexamples.  A failed scoped search is not a lower bound on
all equivalent machines.

\subsection{Proof strength and the internal equivalence obligation}
The published PA196 and PA209 developments distinguish external nonhalting
from what PA proves about the machine.  For PA196 the advertised internal
statement is
\[
 \mathrm{PA}\vdash
 \operatorname{Nonhalt}(M_{196})\longrightarrow\operatorname{Con}(\mathrm{PA}).
\]
The relevant endpoints are the pinned
\href{https://github.com/jamesyc/BusyBeaver/blob/407cdc356f6b586635e71ac92be394255f4bfea4/docs/result.md}{theorem table},
\href{https://github.com/jamesyc/BusyBeaver/blob/407cdc356f6b586635e71ac92be394255f4bfea4/PA196/Independence.lean}{the independence module},
and
\href{https://github.com/jamesyc/BusyBeaver/blob/407cdc356f6b586635e71ac92be394255f4bfea4/PA196/Nonhalting.lean}{the nonhalting module}.
The last file explicitly separates external soundness from arithmetic
unprovability.  The old PA209 article makes the same distinction.

Thus the requested PA-provable biconditional has an additional proof
obligation.  A proof that PA has no proof of either halting or nonhalting,
together with an external soundness proof, does not by itself establish the
missing implication inside PA.  The appropriate additional route is a
primitive recursive syntactic translation from an accepted numerical proof
to a PA refutation, with PA proving the correctness of that translation.
That argument can use finite proof codes and bounded computations; it must
not silently invoke a uniform truth predicate for all PA formulas or PA's
own semantic soundness.  A separate chapter supplies the proposed converse
argument.  The comparison here does not attribute it to an upstream theorem
that states only the displayed implication.

\subsection{The contribution and limits of the ConciseIPRA material}
The pinned
\href{https://github.com/VladimirReshetnikov/ProveIt/tree/6ce5acc401f0c84be90dc5c38dc1487e391122c6/Logic/PrimitiveRecursiveArithmetic/ConciseIPRA}{ConciseIPRA project}
measures complete standalone Python proof checkers by significant tokens,
AST nodes, and normalized source bytes.  Its mathematical criterion is
bi-interpretability with full first-order intuitionistic primitive recursive
arithmetic, with whole numerical domains, actual equality, and identity
composites.  These are substantive requirements, but they are different from
a state census for a binary blank-tape machine with the PA-consistency
property.  The stated bi-interpretability target does not by itself establish
the PA-consistency equivalence needed here.

Several design ideas are directly relevant to minimizing proof searchers:
selecting an operation from the shape of a premise; deriving unrestricted
transport from a one-sided restricted rule; fusing evaluation with discharge;
replacing a general template operation by simpler operations with explicit
proof recovery; and using a syntactic tag to select an introduction mode.
PA196 already uses a related principle when it removes a primitive equality
head but derives its required applications from other heads.

A new use of these ideas would still need two concrete results: a proof
translation preserving the requested PA target, and a literal Turing-machine
compilation with a verified state count.  A smaller rule inventory or shorter
Python checker does not supply either result automatically.  Consequently
this comparison treats the ConciseIPRA proposals as useful design input,
without claiming a state-count improvement from them alone.

\subsection{What the independent audit establishes}
The supplied \texttt{compare\_designs.py} is a standard-library verifier for
all selected transition tables.  It checks row uniqueness, binary actions,
target closure, working-state and halt domains, graph reachability, canonical
graph serialization, exact helper patches, complete predecessor-tag
conditions, and the finite whole-tape local simulations for all eighteen
new helper layouts.  Its checks use explicit failures rather than relying on
Python assertions.  It passes both normally and under \texttt{python -O},
producing byte-identical receipts.

The machine manifest preserves 27 designated 208-state presentations, the
published PA209 and PA196 tables, the two native helper baselines, original
paths, and hashes.  The archive inventory and this chapter keep uncompiled
proposals, failed candidates, scoped search results, local simulation
certificates, and proof-assistant endpoints distinct.  This comparison
chapter establishes no new sub-196 state-count result; any such construction
must be identified and verified independently of the census above.
